\documentclass[11pt]{article}
\usepackage[margin=1.5cm]{geometry}
\usepackage{amsmath,amssymb}
\usepackage{enumitem}
\usepackage{listings}
\usepackage[T1]{fontenc}
\setlength{\columnsep}{0.7cm}
\setlength{\parindent}{0pt}
\setlist{nosep,leftmargin=*}
\lstset{
  basicstyle=\ttfamily\scriptsize,
  columns=fullflexible,
  keepspaces=true,
  aboveskip=2pt,
  belowskip=2pt,
  frame=single,
  breaklines=true
}

\begin{document}
\twocolumn

\title{Chapter 12 Warm-Up: Recursion}
\author{Prof. Jordan C. Hanson}
\date{}

\maketitle

\section{Recursion with a List}
Consider the list \\
\begin{lstlisting}
values = [7, 12, 4, 19, 3, 8, 15, 2, 11, 6, 14, 5]
\end{lstlisting} \vspace{0.25cm}
We want a function that returns \texttt{True} only when \emph{every} number in the list is positive.  First, recall that the following method acts on a list object: \\
\begin{lstlisting}
x = values.pop()
\end{lstlisting}\vspace{0.25cm}
After the statement runs, what value is stored in \texttt{x}, and what has happened to \texttt{values}?  In a recursive algorithm to solve this problem, each function call should examine one number, remove it, and then ask the \emph{same question} about the smaller list. Complete the function: \\
\begin{lstlisting}
def all_positive(values):
    # Base case: no values remain
    if ______________________:
        return ______________

    x = values.pop()

    if ______________________:
        return ______________

    return all_positive(values)
\end{lstlisting}\vspace{0.25cm}

For \\
\begin{lstlisting}
test = [6, 2, 9, 4]
print(all_positive(test))
\end{lstlisting}\vspace{0.25cm}
write the value removed at each recursive call:
\[
\underline{\hspace{0.7cm}}\;\rightarrow\;
\underline{\hspace{0.7cm}}\;\rightarrow\;
\underline{\hspace{0.7cm}}\;\rightarrow\;
\underline{\hspace{0.7cm}}
\]
This exercise is called \textit{tracing the calls.}  What is returned when the list becomes empty? Now predict the result for \\
\begin{lstlisting}
test = [6, 2, -9, 4]
\end{lstlisting}\vspace{0.25cm}
At what point can the recursion stop early?

\section{Recursive Factorial}
For a positive integer \(n\), the factorial is
\[
n! = n(n-1)(n-2)\cdots 2\cdot 1.
\]
For example,
\[
6! = 6\cdot5\cdot4\cdot3\cdot2\cdot1 = 720.
\]

Factorial has a natural recursive definition:
\[
n! = n(n-1)!.
\]
To make the recursion stop, we also need a \textbf{base case}:
\[
0! = 1 \qquad \text{and} \qquad 1! = 1.
\]

Complete the recursive function: \\
\begin{lstlisting}
def factorial(n):
    if ______________________:
        return ______________

    return ___________________
\end{lstlisting}\vspace{0.25cm}

Let us \textit{trace} an example.  Complete the expansion:
\[
\begin{aligned}
6! &= 6(5!)\\
   &= 6\cdot5(\underline{\hspace{0.45cm}}!)\\
   &= 6\cdot5\cdot4(\underline{\hspace{0.45cm}}!)\\
   &= \cdots\\
   &= \underline{\hspace{2.0cm}}.
\end{aligned}
\]

Identify the two parts of recursion:

\begin{enumerate}
\item What is the \textbf{base case} in \texttt{factorial()}?
\item What is the \textbf{recursive case}?
\item In each recursive call, how does the problem become smaller?
\end{enumerate}

\section{Challenge Problem}

What should happen if a user calls \texttt{factorial(-3)}? Add one conditional that prevents the function from recursing forever on a negative input.

\vfill
\small\textbf{Key idea:} A recursive function must have a stopping case
and a recursive step that moves toward that stopping case.

\end{document}
